\section{Conclusions and Final Architecture Selection}
\label{sec:conclusions}

\paragraph{What the measurements support.}
On the host of \cref{tab:host}, over the \fixed{100}-book corpus and the
\fixed{30}-query workload, the tournament measured the following:

\begin{itemize}
    \item The four implementations are behaviourally identical on this
    workload (byte-identical canonical outputs, \cref{sec:equivalence}), so
    the cross-language differences below compare the same work.
    \item On write throughput and datalake lookup, the measured spread
    between implementations is \pending{throughput/lookup summary}.
    \item On index build time and query latency, \pending{build/query
    summary}; on incremental updates over the measured base sizes,
    \pending{update summary}.
    \item On-disk footprint is dominated by the corpus itself; the storage
    overhead of the three index structures over the same corpus is
    \pending{index storage summary} (\cref{sec:index-memory}).
    \item Recovery after a documented interruption resumed with
    \pending{0} duplicated and \pending{0} lost books in every layout
    (\cref{tab:recovery}).
\end{itemize}

\paragraph{Selection for Stage 2.}
Within the tested corpus, workload and host, the group selects
\textbf{Java with the MongoDB index structure} for Stage 2. The reasoning is
scoped to the evidence: MongoDB showed \pending{measured update and space
behaviour} at the measured sizes, its server-side footprint was measured
separately (\texttt{collStats}), and the Java port's end-to-end numbers were
supported by JMH attribution. This is a choice \emph{for the Stage 2
workload}, not a claim that Java is universally faster or that MongoDB is
the only scalable storage: at the measured scale, the file-based structures
were competitive on several axes, and behaviours that matter at larger
scale --- sharding, replication, concurrency, multi-gigabyte indexes ---
were \emph{not measured here} and remain open questions
(\cref{sec:limitations}).

\paragraph{Language-wide claims.}
Nothing in this report supports a statement of the form ``language X is
faster than language Y''. The measured objects are four implementations,
written by this team, of one pipeline, on one corpus, on one machine, in one
continuous run. Statements such as ``Go's allocator showed the lowest RSS
growth in this workload'' are observations; ``Go uses less memory'' is a
hypothesis this experiment cannot test. Where the data was ambiguous, the
report says so rather than picking a narrative.

\subsection{Future Improvements}
\label{sec:future}

\begin{itemize}
    \item Repeat the tournament on bare metal and on Linux natively to
    separate runtime effects from WSL2's virtualisation.
    \item Extend the corpus by an order of magnitude to observe where the
    measured curves bend (the current 10{,}000-book batch reuses a
    100-book corpus, \cref{sec:workload}).
    \item Measure concurrency (multiple query clients) for the Stage 2
    serving layer, where it will actually be used.
    \item Add a memory-profiled MongoDB run that accounts the server's RSS
    alongside the client's, for a true capacity figure.
\end{itemize}
